¿Qué acciones haces para protegerte de ciberataques?

\textbf{Respuesta:}
Previo a cursar lo poco que he cursado de la materia, no he aplicado muchas medidas para protegerme de ciberataques, sin embargo, algunas de las que utilizo son:

\begin{itemizeColor}
    \item \textbf{Autenticación multifactor (MFA/2FA):} Habilitar el doble factor de autenticación en todas las cuentas posibles, priorizando el uso de aplicaciones manejadoras de contraseñas seguras, y evitando la verificación por SMS debido al riesgo de suplantación de SIM.
    \item \textbf{Gestión de contraseñas seguras y únicas:} Emplear un gestor de contraseñas para crear y almacenar contraseñas aleatorias y complejas, erradicando por completo la reutilización de claves ante posibles filtraciones de datos.
    \item \textbf{Gestión de parches y actualizaciones periódicas:} Mantener actualizados el sistema operativo, los navegadores web y las aplicaciones para mitigar vulnerabilidades, realizo actualizaciones al menos uno o dos días después de salir la última versión, también para esperar a que salga información con respecto a la nueva actualización y asegurar que no se me exponga a vulnerabilidades de la actualización que puedan ser más graves.
    \item \textbf{Navegación precavida y verificación de enlaces:} Soy desconfiado frente a correos y mensajes con carácter urgente o sospechoso, me tomo mi tiempo para analizar el dominio remitente y la URL antes de abrir enlaces o descargar archivos adjuntos.
    \item \textbf{Copias de seguridad periódicas (estrategia 3-2-1):} Mantener respaldos frecuentes de los datos más importantes, incluyendo copias en almacenamiento local desconectado de la red (\textit{offline}) para proteger la información contra ataques de secuestro de datos.
\end{itemizeColor}

